% !TEX program = pdflatex
% Hypothèses de logique métier du pipeline graph_pipeline_prototype.ipynb
% Compilation : pdflatex pipeline_hypotheses.tex (2 passes pour les références)

\documentclass[11pt,a4paper]{article}

\usepackage[utf8]{inputenc}
\usepackage[T1]{fontenc}
\usepackage[french]{babel}
\usepackage{lmodern}
\usepackage[margin=2.2cm]{geometry}
\usepackage{amsmath,amssymb,amsfonts,mathtools}
\usepackage{booktabs}
\usepackage{array}
\usepackage{longtable}
\usepackage{listings}
\usepackage{xcolor}
\usepackage{enumitem}
\usepackage[hidelinks]{hyperref}
\usepackage{caption}
\usepackage{tikz}
\usetikzlibrary{positioning,arrows.meta,fit,calc,backgrounds}

% Désactive les espacements automatiques du français devant ; : ! ?
\shorthandoff*{;:!?}

\definecolor{codebg}{RGB}{247,247,250}
\definecolor{codekw}{RGB}{20,80,160}
\definecolor{codecm}{RGB}{110,120,130}
\definecolor{codestr}{RGB}{150,60,60}
\definecolor{hypcol}{RGB}{150,40,40}
\definecolor{srccol}{RGB}{90,90,95}

\lstdefinestyle{conf}{
  backgroundcolor=\color{codebg},
  basicstyle=\ttfamily\scriptsize,
  keywordstyle=\color{codekw},
  commentstyle=\color{codecm}\itshape,
  stringstyle=\color{codestr},
  breaklines=true,
  columns=fullflexible,
  frame=single,
  framesep=4pt,
  rulecolor=\color{codecm!40},
  numbers=none,
  showstringspaces=false,
  tabsize=2,
}
\lstset{style=conf}

\newcommand{\code}[1]{\texttt{\small #1}}
\newcommand{\hyp}[1]{\textcolor{hypcol}{\textbf{\small [#1]}}}
\newcommand{\cell}[1]{\textbf{Cellule #1}}
\newcommand{\htg}[1]{{\sffamily\tiny\bfseries\color{hypcol}#1}}

\tikzset{
  tbl/.style={draw=srccol!60, fill=srccol!8, rounded corners=1pt,
              font=\ttfamily\tiny, align=center, inner sep=3pt, text width=2.9cm},
  fn/.style={draw=codekw!75, fill=codekw!8, rounded corners=3pt, thick,
             font=\ttfamily\scriptsize, align=center, inner sep=4pt, text width=3.1cm},
  op/.style={draw=srccol!70, fill=white, rounded corners=1pt,
             font=\sffamily\tiny, align=center, inner sep=3pt, text width=3.1cm},
  art/.style={draw=codekw!40, fill=codekw!4, rounded corners=1pt,
              font=\ttfamily\tiny, align=center, inner sep=3pt, text width=3.1cm},
  htag/.style={font=\sffamily\tiny\bfseries, text=hypcol},
  ar/.style={-{Stealth[length=5pt,width=4pt]}, draw=srccol!85, thick},
  arw/.style={-{Stealth[length=5pt,width=4pt]}, draw=srccol!60, dashed},
  lbl/.style={font=\sffamily\tiny, text=srccol, fill=white, inner sep=1.5pt},
}

\title{\vspace{-1.6cm}\textbf{Pipeline de construction du graphe de co-vues}\\[0.25cm]
\large Enchaînement des fonctions, hypothèses de logique métier\\
et sources de variabilité entre deux exécutions}

\author{Neil Braun\\[0.15cm]
\small Mirakl --- Data Science / Target2Sell}

\date{\small Août 2026 --- source : \code{graph\_pipeline\_prototype.ipynb}}

\begin{document}

\maketitle

\noindent
Ce document part du seul notebook \code{graph\_pipeline\_prototype.ipynb}. Il isole
d'abord les deux chaînes de fonctions qu'il contient --- celle qui construit le
\textbf{graphe} et celle qui construit les \textbf{arêtes d'entraînement} --- puis
retrace le notebook cellule par cellule en renvoyant, à chaque étape, aux hypothèses
de logique métier numérotées \hyp{H1}\dots\hyp{H34} (sections~\ref{sec:hyp-a}
et~\ref{sec:hyp-b}). La section~\ref{sec:exclusions} recense tout ce qui peut être
silencieusement exclu (\hyp{E1}\dots\hyp{E19}) et la section~\ref{sec:variabilite}
explique pourquoi deux exécutions du notebook, même sur la même fenêtre de dates,
ne donnent pas le même graphe (\hyp{V1}\dots\hyp{V14}).

\medskip
\noindent
Toute affirmation est tirée du code du notebook. Les rares points qui dépendent de
l'environnement (version de Spark, rétention des tables) sont signalés comme
\emph{à vérifier} et non présentés comme relevés.

\tableofcontents

%======================================================================
\section{Les deux chaînes de fonctions}

Le notebook contient deux pipelines distincts qui ne se rejoignent qu'à la toute
fin, dans l'objet \code{Data} de PyTorch Geometric :

\begin{description}[leftmargin=1.1cm,style=nextline,font=\normalfont\itshape]
  \item[Chaîne A --- le graphe.] Tracking de navigation $\rightarrow$ sessions
        $\rightarrow$ paires de produits co-vus $\rightarrow$ \code{edge\_index}
        d'entraînement + matrice de features \code{x}. C'est le graphe sur lequel
        circule le \emph{message passing}.
  \item[Chaîne B --- les arêtes d'entraînement.] Logs publicitaires
        (\code{ads\_adlog\_fct}) rapprochés du tracking $\rightarrow$ exécutions de
        recommandation $\rightarrow$ paires (trigger, candidat) étiquetées
        positif/négatif. C'est ce que la tête de prédiction apprend à classer, et ce
        sur quoi le MRR est calculé.
\end{description}

\noindent
Ces deux chaînes n'ont pas la même sémantique d'arête : la chaîne~A produit
« ces deux produits ont été vus dans la même session », la chaîne~B produit
« ce produit a été recommandé à la suite de ce trigger, et vu / non vu »
\hyp{H34}.

\subsection{Chaîne A --- construction du graphe}

\begin{center}
\begin{tikzpicture}[font=\small]

% --- sources
\node[tbl] (s1) at (-6.4,1.5) {t2s\_gold\_customer\\artemis\_insights\_customer\\artemis\_publisher};
\node[tbl] (s2) at (-2.9,1.5) {t2s\_mongo\_product\\\_0\_current};
\node[tbl] (s3) at (0.6,1.5) {tracking\_product\_display\\\_page\_event\_fct\\t2s\_merged\_user};
\node[tbl] (s4) at (4.1,1.5) {product\_embeddings};

% --- fonctions d'extraction
\node[fn] (f1) at (-6.4,0.0) {get\_customer\\[1pt]\htg{H1}};
\node[fn] (f2) at (-2.9,0.0) {get\_products\\[1pt]\htg{H2 H3 H4}};
\node[fn] (f3) at (0.6,0.0) {get\_views\\[1pt]\htg{H6 H7 H8 H9}};
\node[fn] (f4) at (4.1,0.0) {get\_embeddings\\[1pt]\htg{H5}};

\draw[ar] (s1) -- (f1);
\draw[ar] (s2) -- (f2);
\draw[ar] (s3) -- (f3);
\draw[ar] (s4) -- (f4);

% id resolution
\draw[arw] (f1.south) -- ++(0,-0.45) -| node[lbl,pos=0.25]{db\_names, customerId} (f2.south);
\draw[arw] (f1.south) -- ++(0,-0.45) -| (f3.south);

% --- jointure catalogue
\node[op] (j1) at (-1.15,-1.55) {jointure \code{inner} sur \code{internalId}\\(vues $\cap$ catalogue) \htg{H11}};
\draw[ar] (f2.south) |- ++(0,-1.1) -- (j1.north west);
\draw[ar] (f3.south) |- ++(0,-1.1) -- (j1.north east);

\node[fn] (f5) at (-1.15,-2.85) {sessionize\_views\\[1pt]\htg{H10}};
\node[fn] (f6) at (-1.15,-4.15) {split\_sessions\\[1pt]\htg{H17}};
\node[art] (dsplit) at (-1.15,-5.25) {\textbf{df\_split} (\code{cache})\\userId, session\_id, internalId,\\emitted, session\_start, split};

\draw[ar] (j1) -- (f5);
\draw[ar] (f5) -- (f6);
\draw[ar] (f6) -- (dsplit);

% --- build co-view edges x3
\node[fn] (b1) at (-5.0,-6.85) {build\_co\_view\_edges\\\scriptsize(split = train)};
\node[fn] (b2) at (-1.15,-6.85) {build\_co\_view\_edges\\\scriptsize(split = val)};
\node[fn] (b3) at (2.7,-6.85) {build\_co\_view\_edges\\\scriptsize(split = test)};
\node[htag] at (-1.15,-7.75) {H12 H13 H14 H15 H16};

\draw[ar] (dsplit.south) -- ++(0,-0.35) -| (b1.north);
\draw[ar] (dsplit.south) -- ++(0,-0.35) -| (b2.north);
\draw[ar] (dsplit.south) -- ++(0,-0.35) -| (b3.north);

\node[fn] (r1) at (0.8,-8.6) {remove\_train\_edges\\\scriptsize(\code{left\_anti} sur les paires train)\\[1pt]\htg{H18}};
\draw[ar] (b2.south) -- ++(0,-0.5) -| (r1.north west);
\draw[ar] (b3.south) -- ++(0,-0.5) -| (r1.north east);
\draw[arw] (b1.south) -- ++(0,-0.9) -| node[lbl,pos=0.2]{paires à retirer} (r1.west);

\node[art] (nodes) at (-1.15,-9.9) {\textbf{df\_graph\_nodes} = union des\\extrémités train $\cup$ val $\cup$ test \htg{H19}};
\draw[ar] (b1.south) -- ++(0,-1.9) -| (nodes.west);
\draw[ar] (r1.south) -- ++(0,-0.55) -| (nodes.east);

\node[fn] (m1) at (-1.15,-11.2) {build\_node\_mapping\\\scriptsize$\rightarrow$ \code{node2idx} (tri des internalId)\\[1pt]\htg{H20}};
\draw[ar] (nodes) -- (m1);

\node[fn] (rm) at (-4.3,-12.7) {remap\_edges};
\node[fn] (emb) at (2.0,-12.7) {\code{node\_features} $\rightarrow$ \code{x}\\\scriptsize(0 si pas d'embedding) \htg{H5}};
\draw[ar] (m1.south) -- ++(0,-0.35) -| (rm.north);
\draw[ar] (m1.south) -- ++(0,-0.35) -| (emb.north);
\draw[ar] (f4.east) -- (6.1,0) -- node[lbl,pos=0.55,rotate=90]{embeddings $\cap$ nœuds du graphe} (6.1,-12.7) -- (emb.east);

\node[fn] (et) at (-4.3,-13.9) {edges\_to\_tensors\\\scriptsize(\code{symmetric=True}) \htg{H16}};
\draw[ar] (rm) -- (et);

\node[fn] (pyg) at (-1.15,-15.3) {build\_pyg\_data $\rightarrow$ \code{Data}};
\draw[ar] (et.south) -- ++(0,-0.45) -| (pyg.north);
\draw[ar] (emb.south) -- ++(0,-0.45) -| (pyg.north);

\node[art] (ds) at (-1.15,-16.4) {\textbf{CoViewDataset} $\rightarrow$ \code{data.pt} + \code{node2idx.json}};
\draw[ar] (pyg) -- (ds);

\end{tikzpicture}
\end{center}

\begin{center}
\parbox{0.92\textwidth}{\footnotesize\itshape
Figure 1 --- Chaîne A. Les cadres bleus sont des fonctions du notebook, les cadres
gris des tables Unity Catalog, les cadres clairs des artefacts intermédiaires. Les
étiquettes rouges renvoient aux hypothèses du
tableau~\ref{tab:hyp-a}. \code{df\_val\_edges} et \code{df\_test\_edges} ne servent
qu'à élargir le jeu de nœuds et à produire les compteurs : ils ne sont pas passés à
\code{build\_pyg\_data} (voir \hyp{H19}).}
\end{center}

\clearpage
\subsection{Chaîne B --- construction des arêtes d'entraînement}

\begin{center}
\begin{tikzpicture}[font=\small]

\node[tbl] (t1) at (-5.4,1.4) {ads\_adlog\_fct};
\node[tbl] (t2) at (-1.8,1.4) {tracking\_product\_display\\\_page\_event\_fct};
\node[tbl] (t3) at (1.8,1.4) {t2s\_merged\_user};

\node[fn] (rec) at (-1.8,-0.4) {recommended\_products\_for\_t2s\_user\\\scriptsize jointure $\pm 3$\,s sur (productId, cookie, customerId)\\[1pt]\htg{H21 H22}};
\draw[ar] (t1.south) -- ++(0,-0.6) -| (rec.north west);
\draw[ar] (t2) -- (rec.north);
\draw[ar] (t3.south) -- ++(0,-0.6) -| (rec.north east);

\node[op] (j2) at (-1.8,-2.1) {jointure \code{inner} avec \code{df\_split}\\sur (internalId, emitted, userId) \htg{H23}};
\draw[ar] (rec) -- (j2);
\draw[arw] (-5.6,-2.1) -- node[lbl,pos=0.5]{df\_split (chaîne A)} (j2.west);

\node[art] (drec) at (-1.8,-3.4) {\textbf{df\_recommended\_products\_t2s}\\\textbf{\_with\_sessions\_indexes} (\code{cache})};
\draw[ar] (j2) -- (drec);

% --- bloc positive_negative_edge_construction
\node[fn] (ex) at (-4.6,-5.3) {explode(placements)\\+ slice top-12 + posexplode\\[1pt]\htg{H24 H25}};
\node[art] (dre) at (-4.6,-6.7) {df\_recommended\_exploded\\\scriptsize (trigger, candidat, rang)};
\draw[ar] (drec.south) -- ++(0,-0.5) -| (ex.north);
\draw[ar] (ex) -- (dre);

\node[fn] (pos) at (-7.0,-8.5) {positifs\\\scriptsize candidat $\in$ produits de la session\\[1pt]\htg{H26}};
\node[fn] (neg) at (-2.6,-8.5) {négatifs\\\scriptsize $\lnot$session $\land$ $\lnot$arête train $\land$ $\lnot$positif\\[1pt]\htg{H27 H28}};
\draw[ar] (dre.south) -- ++(0,-0.45) -| (pos.north);
\draw[ar] (dre.south) -- ++(0,-0.45) -| (neg.north);
\draw[arw] (1.1,-8.5) -- node[lbl,pos=0.45]{df\_train\_edges} (neg.east);

\node[fn] (mrr) at (-7.0,-10.2) {MRR production\\\scriptsize matched-only / coverage-adjusted};
\draw[ar] (pos) -- (mrr);

\node[fn] (trunc) at (-2.6,-10.2) {troncature \code{rank} $\le$ \code{deepest\_pos\_rank}\\\scriptsize + \code{sample(fraction=0.0)}\\[1pt]\htg{H29}};
\draw[ar] (neg) -- (trunc);

\node[fn] (ec) at (-4.6,-11.9) {\code{exec2code} sur toutes les exécutions\\\scriptsize (tri $\rightarrow$ déterminisme) \htg{H31}};
\draw[ar] (drec.east) -- (1.6,-3.4) -- node[lbl,pos=0.6,rotate=90]{toutes les executionId du split} (1.6,-11.9) -- (ec.east);

\node[fn] (mapn) at (-4.6,-13.4) {\code{map(node2idx)} + \code{dropna}\\[1pt]\htg{H30}};
\draw[ar] (mrr.south) -- ++(0,-2.4) -| (mapn.west);
\draw[ar] (trunc.south) -- ++(0,-2.4) -| (mapn.east);
\draw[ar] (ec) -- (mapn);

\node[art] (tens) at (-4.6,-14.7) {\code{pos\_edge\_index}, \code{neg\_edge\_index}\\\scriptsize $[3, N]$ = (exec\_code, trigger, produit)};
\draw[ar] (mapn) -- (tens);

\begin{scope}[on background layer]
  \node[draw=hypcol!45, dashed, rounded corners=4pt, fit=(ex)(pos)(neg)(mrr)(trunc)(ec)(mapn)(tens),
        inner sep=7pt, label={[font=\sffamily\scriptsize\itshape, text=hypcol]above right:positive\_negative\_edge\_construction (appelée pour train, val, test)}] {};
\end{scope}

\end{tikzpicture}
\end{center}

\begin{center}
\parbox{0.92\textwidth}{\footnotesize\itshape
Figure 2 --- Chaîne B. Le cadre pointillé délimite l'unique fonction
\code{positive\_negative\_edge\_construction}, appelée trois fois (train, val,
test), toujours avec \code{df\_train\_edges} comme référence de négatifs.}
\end{center}

\subsection{Chaîne B$'$ --- tables de référence production (split \emph{val})}

\begin{center}
\begin{tikzpicture}[font=\small]

\node[art] (drec2) at (0,0.8) {\textbf{df\_recommended\_products\_t2s\_with\_sessions\_indexes}, \code{split = val}};

\node[fn] (vt) at (-5.0,-0.9) {val\_triggers\\\scriptsize\code{distinct} (user, session, exec, trigger)};
\node[fn] (vp) at (3.5,-0.9) {val\_placements\_prototype\\\scriptsize posexplode(placements)\\+ \code{dropDuplicates(exec, pos)} \htg{H25}};
\draw[ar] (drec2.south) -- ++(0,-0.35) -| (vt.north);
\draw[ar] (drec2.south) -- ++(0,-0.35) -| (vp.north);

\node[fn] (vsp) at (-5.0,-2.7) {val\_session\_products\\\scriptsize \code{collect\_list} des produits de la session};
\draw[ar] (vt) -- (vsp);
\draw[arw] (-8.0,-2.7) -- node[lbl,pos=0.5]{df\_split} (vsp.west);

\node[art] (sr) at (-5.0,-4.4) {\textbf{sessions\_raw\_val\_prototype}\\\scriptsize $\rightarrow$ parquet S3};
\draw[ar] (vsp) -- (sr);
\draw[arw] (-1.6,-4.4) -- node[lbl,pos=0.5]{val\_exec2code} (sr.east);

\node[fn] (pd) at (1.2,-2.7) {prod\_products\_by\_trigger\\\scriptsize \code{cands\_display} (ordre d'affichage)};
\node[fn] (pr) at (5.8,-2.7) {prod\_products\_by\_trigger\\\scriptsize \code{cands\_relevance} (relevanceScore)};
\draw[ar] (vp.south) -- ++(0,-0.5) -| (pd.north);
\draw[ar] (vp.south) -- ++(0,-0.5) -| (pr.north);

\node[fn] (apn) at (3.5,-4.4) {add\_pos\_neg\\\scriptsize positif = produit renvoyé $\in$ session\\[1pt]\htg{H32}};
\draw[ar] (pd.south) -- ++(0,-0.5) -| (apn.north west);
\draw[ar] (pr.south) -- ++(0,-0.5) -| (apn.north east);

\node[fn] (dd) at (3.5,-6.0) {\code{dropDuplicates(["exec\_code"])}\\[1pt]\htg{H33}};
\draw[ar] (sr.south) -- ++(0,-0.5) -| (dd.west);
\draw[ar] (apn) -- (dd);

\node[art] (pv) at (3.5,-7.3) {\textbf{prod\_results\_val\_prototype}\\\scriptsize $\rightarrow$ parquet S3};
\draw[ar] (dd) -- (pv);

\end{tikzpicture}
\end{center}

\begin{center}
\parbox{0.92\textwidth}{\footnotesize\itshape
Figure 3 --- Chaîne B$'$. Ces deux parquets ne servent pas à l'entraînement : ils
donnent la baseline production sur \emph{val}, en deux variantes d'ordre des 12
candidats. Aucun filtrage sur les arêtes du graphe n'y est appliqué, contrairement à
la chaîne~B (\hyp{H27} vs \hyp{H32}).}
\end{center}

%======================================================================
\clearpage
\section{Parcours du notebook, cellule par cellule}

Chaque étape est résumée en une ligne, suivie des hypothèses qu'elle pose.
Les libellés complets sont dans les tableaux~\ref{tab:hyp-a} et~\ref{tab:hyp-b}.

\subsection*{Cellules 0--2 --- session Spark et définition des fonctions}
Aucune exécution de logique métier : \code{DatabricksSession}, installation de
\code{torch\_geometric}, puis définition des sept fonctions de la chaîne~A.

\subsection*{\cell{3} --- résolution du client et fenêtre temporelle}
\begin{itemize}[nosep,leftmargin=*]
  \item \code{get\_customer("maisons-du-monde")} puis \code{collect()} :
        les \code{customerId}, \code{publisherId} et \code{db\_name} sont ramenés au
        driver sous forme de listes. \hyp{H1}
  \item \code{get\_products} et \code{get\_embeddings} sont câblés sur ces
        \code{db\_names}. \hyp{H2} \hyp{H3} \hyp{H4} \hyp{H5}
  \item Fenêtre : \code{start\_date = 2026-01-01}, \code{end\_date = 2026-05-01},
        donc \code{cutoff\_train = 2026-04-01} et \code{cutoff\_val = 2026-04-16}.
        Train $\approx$ 90 jours, val 15 jours, test 15 jours. \hyp{H9} \hyp{H17}
  \item \code{get\_views} charge le tracking brut réconcilié par
        \code{t2s\_merged\_user}. \hyp{H6} \hyp{H7} \hyp{H8}
\end{itemize}

\subsection*{\cell{4} --- sessions, splits et arêtes de co-vue}
\begin{itemize}[nosep,leftmargin=*]
  \item \code{df\_views.join(df\_products, "internalId", "inner")} : le filtre
        catalogue est appliqué \emph{avant} la sessionisation, il peut donc scinder
        des sessions. \hyp{H11}
  \item \code{sessionize\_views} : coupure à 1800\,s d'inactivité, fenêtre
        partitionnée par \code{userId} et ordonnée par
        \code{(emitted, internalId)}. \hyp{H10}
  \item \code{split\_sessions} : affectation train / val / test par
        \code{session\_start}, puis \code{cache()}. \hyp{H17}
  \item \code{build\_co\_view\_edges} sur chacun des trois sous-ensembles, puis
        \code{remove\_train\_edges} sur val et test.
        \hyp{H12} \hyp{H13} \hyp{H14} \hyp{H15} \hyp{H18}
  \item \code{count\_nodes} sur l'union des trois jeux d'arêtes : le compteur affiché
        inclut donc des nœuds qui n'auront aucune arête d'entraînement. \hyp{H19}
\end{itemize}

\subsection*{\cell{5} --- passage en pandas}
\begin{itemize}[nosep,leftmargin=*]
  \item \code{df\_graph\_nodes} = union des six colonnes d'extrémités
        (train, val, test $\times$ A, B), \code{distinct}. \hyp{H19}
  \item \code{toPandas()} sur les trois jeux d'arêtes : tout le graphe passe sur le
        driver. Les embeddings sont restreints par jointure \code{inner} aux nœuds du
        graphe.
\end{itemize}

\subsection*{Cellules 6--7 --- indexation des nœuds}
\begin{itemize}[nosep,leftmargin=*]
  \item \code{build\_node\_mapping} concatène les extrémités des trois jeux,
        \code{unique()} puis \code{sort()} : l'index est l'ordre croissant des
        \code{internalId}. \hyp{H20}
  \item \code{remap\_edges} remplace les \code{internalId} par ces index. Un
        identifiant absent du mapping devient \code{NaN} --- impossible ici, le
        mapping étant construit sur ces mêmes tables.
\end{itemize}

\subsection*{\cell{8} --- matrice de features}
\begin{itemize}[nosep,leftmargin=*]
  \item \code{embedding\_dim} est lu sur la \emph{première} ligne de
        \code{emb\_pandas} : dimension supposée homogène, garantie ici par le filtre
        sur \code{last\_model\_id}. \hyp{H5}
  \item \code{node\_features} est initialisée à zéro puis remplie par affectation
        indexée. Un nœud sans embedding garde un vecteur nul, et le compteur
        \code{Nodes without embedding} le mesure. \hyp{H5}
  \item Si \code{emb\_pandas} contenait deux lignes pour un même \code{internalId},
        la dernière écrase la précédente, dans l'ordre de collecte de Spark ---
        non déterministe. \hyp{V10}
\end{itemize}

\subsection*{\cell{9} --- tensorisation}
\begin{itemize}[nosep,leftmargin=*]
  \item \code{edges\_to\_tensors(symmetric=True)} duplique chaque paire dans les deux
        sens : $E$ double. \hyp{H16}
  \item \code{val\_edge\_index} et \code{test\_edge\_index} sont calculés ici mais ne
        sont \textbf{pas} passés à \code{build\_pyg\_data} (\cell{25}). \hyp{H19}
\end{itemize}

\subsection*{Cellules 10--11 --- rapprochement avec les logs publicitaires}
\begin{itemize}[nosep,leftmargin=*]
  \item \code{recommended\_products\_for\_t2s\_user} : jointure \code{inner} entre
        \code{ads\_adlog\_fct} et le tracking sur
        (\code{productId} = \code{fwProductId}, \code{userCookie} =
        \code{user.cookie}, \code{customerId}) dans une fenêtre de $\pm 3$ secondes
        autour de \code{logInstant}, restreinte à \code{pageType = "PRODUCT"}.
        \hyp{H21} \hyp{H22}
  \item Jointure avec \code{df\_split} sur le triplet exact
        (\code{internalId}, \code{emitted}, \code{userId}) : chaque exécution hérite
        de son \code{session\_id} et de son \code{split}. \hyp{H23}
\end{itemize}

\subsection*{\cell{12} --- positifs, négatifs et MRR production}
\begin{itemize}[nosep,leftmargin=*]
  \item Jointure trigger $\times$ produits de la session, sur
        (\code{userId}, \code{session\_id}).
  \item \code{explode(sponsoredProductPlacementExecutions)}, puis
        \code{slice(relevantProducts, 1, 12)} et \code{posexplode} pour obtenir un
        rang 1-based. \hyp{H24} \hyp{H25}
  \item \textbf{Positif} : \code{internalId\_session = internalId\_candidate}, sans
        aucune contrainte d'ordre temporel entre la vue et la recommandation.
        \hyp{H26}
  \item MRR production : rang minimal d'un positif par \code{executionId}, en deux
        conventions (\emph{matched-only}, comparable à \code{MRR\_trigger} du modèle,
        et \emph{coverage-adjusted}). Le dénominateur ne compte que les exécutions
        présentes dans \code{df\_recommended\_exploded}. \hyp{E14}
  \item \textbf{Négatif} : candidat non vu dans la session, non présent dans
        \code{df\_train\_edges} (testé dans les deux sens), et non déjà positif.
        Les positifs, eux, ne sont pas filtrés sur les arêtes existantes.
        \hyp{H27} \hyp{H28}
  \item Réduction du volume de négatifs : on ne garde que les rangs
        $\le$ \code{deepest\_pos\_rank} de l'exécution, et les exécutions sans aucun
        positif sont échantillonnées à \code{fraction=0.0}, donc totalement
        écartées. \hyp{H29}
  \item \code{exec2code} est construit sur l'union positifs $\cup$ négatifs
        $\cup$ toutes les exécutions du split, puis \emph{trié} --- le commentaire du
        code indique explicitement que sans ce tri le code dépendrait de l'ordre de
        collecte de Spark. \hyp{H31}
  \item \code{map(node2idx)} puis \code{dropna} : les candidats absents du graphe
        disparaissent des positifs comme des négatifs. \hyp{H30}
\end{itemize}

\subsection*{Cellules 13--16 --- \code{sessions\_raw\_val\_prototype}}
\begin{itemize}[nosep,leftmargin=*]
  \item \code{val\_triggers} : triggers distincts du split val.
  \item \code{val\_session\_products} : pour chaque trigger, la liste
        (\code{internal\_id}, \code{emitted}) de tous les produits de sa session ---
        y compris ceux vus \emph{après} le trigger. \hyp{H26}
  \item Jointure \code{inner} sur \code{exec\_code} et contrôle de couverture
        (\code{n\_triggers - n\_sessions} doit valoir 0). Écriture en parquet sur S3.
\end{itemize}

\subsection*{Cellules 17--21 --- \code{prod\_results\_val\_prototype}}
\begin{itemize}[nosep,leftmargin=*]
  \item \code{val\_placements\_prototype} : un placement identifié par
        (\code{executionId}, \code{placement\_pos}), dédoublonné car la jointure
        session duplique les lignes adlog. \hyp{H25}
  \item Deux variantes du même tableau \code{relevantProducts} :
        \code{cands\_display} (ordre du tableau) et \code{cands\_relevance}
        (tri par \code{relevanceScore} décroissant, avec \code{internalId} comme
        départage). \hyp{H24}
  \item \code{prod\_products\_by\_trigger} fusionne les placements d'un trigger, prend
        le rang minimal par produit, puis \code{row\_number} et coupe à 12.
  \item \code{add\_pos\_neg} : positif = produit renvoyé appartenant à
        \code{session\_ids}, sans aucun filtrage sur les arêtes du graphe. \hyp{H32}
  \item \code{dropDuplicates(["exec\_code"])} : une exécution est réduite à un seul
        trigger, choisi arbitrairement. \hyp{H33}
\end{itemize}

\subsection*{Cellules 22--23 --- diagnostic des deux variantes}
Comparaison des ensembles \code{display} / \code{relevance}
(\code{pct\_identical\_sets}, MRR dans les deux conventions). Si
\code{pct\_identical\_sets} est proche de 1, les deux variantes sont équivalentes et
la comparaison en aval est sans objet --- c'est un test de l'hypothèse \hyp{H24}.

\subsection*{Cellules 24--29 --- assemblage et publication}
\begin{itemize}[nosep,leftmargin=*]
  \item \code{build\_pyg\_data} place les co-vues train dans \code{edge\_index} et les
        paires ads dans \code{train\_pos\_edge\_index} /
        \code{train\_neg\_edge\_index}. \hyp{H34}
  \item \code{CoViewDataset} sérialise \code{data.pt} et \code{node2idx.json}, puis
        \code{data\_prototype.pt}, \code{node2idx\_prototype.json} et
        \code{exec\_mappings\_prototype.json} sont poussés sur S3 sous des chemins
        suffixés \code{\_prototype} --- précisément parce que les \code{exec\_code}
        diffèrent du pipeline principal. \hyp{H31} \hyp{V13}
\end{itemize}

%======================================================================
\clearpage
\section{Hypothèses de la chaîne A --- le graphe}
\label{sec:hyp-a}

{\small
\begin{longtable}{@{}l p{2.45cm} p{5.9cm} p{4.5cm}@{}}
\toprule
\textbf{Id} & \textbf{Étape} & \textbf{Hypothèse métier} & \textbf{Si elle est fausse} \\
\midrule
\endfirsthead
\toprule
\textbf{Id} & \textbf{Étape} & \textbf{Hypothèse métier} & \textbf{Si elle est fausse} \\
\midrule
\endhead
\bottomrule
\endfoot

H1 & \code{get\_customer} & Le \code{shortName} désigne un client unique, qui possède un \emph{publisher} dans le namespace \code{artemis-prod}. & La jointure \code{inner} rend le client invisible, ou \code{collect()} renvoie plusieurs \code{db\_name} et le graphe mélange deux catalogues. \\
\addlinespace
H2 & \code{get\_products} & Le catalogue \emph{courant} vaut pour toute la fenêtre historique : un produit actif aujourd'hui l'était déjà lors des vues de janvier. & Les produits désactivés depuis sont retirés du graphe alors que leurs vues existaient : le graphe rétrécit à mesure que l'on relance tard. \\
\addlinespace
H3 & \code{get\_products} & Seuls les produits maîtres (\code{parentId = -1}) portent le signal ; le tracking a déjà ramené les vues de déclinaisons au maître. & Le signal des déclinaisons est perdu, ou au contraire agrégé alors que les déclinaisons se comportent différemment. \\
\addlinespace
H4 & \code{get\_products} & \code{isActive} $\land$ \code{isOkInStream} = « produit recommandable » ; les autres n'ont pas à être des nœuds. & Des produits recommandés en production sont absents du graphe, donc non scorables (voir \hyp{H30}). \\
\addlinespace
H5 & \code{get\_embeddings}, \cell{8} & Un seul modèle d'embedding (\code{last\_model\_id} figé), une ligne par produit, dimension homogène ; embedding absent $\rightarrow$ vecteur nul. & Doublons $\rightarrow$ écrasement non déterministe ; le vecteur nul est traité par le modèle comme un point du même espace, pas comme un « inconnu ». \\
\addlinespace
H6 & \code{get\_views} & Un événement \code{product\_display\_page\_event} = une vue produit. Pas de pondération par durée ou scroll, pas de déduplication des rafraîchissements. & Un rechargement de page compte comme une vue supplémentaire ; les vues de bots gonflent les sessions (voir \hyp{H14}). \\
\addlinespace
H7 & \code{get\_views} & \code{t2s\_merged\_user} est fonctionnel (un \code{sourceId} $\rightarrow$ un \code{masterId}) et l'identité réconciliée est la bonne granularité de session. & Plusieurs \code{masterId} pour un \code{sourceId} dupliquent les lignes de vues et \emph{doublent} les poids d'arêtes. Fusionner les appareils enchaîne des navigations parallèles dans une même session. \\
\addlinespace
H8 & \code{get\_views} & Trafic non identifié (\code{user.internalId} null) et vue non résolue en produit maître ne portent pas de signal exploitable. & Une part importante du trafic --- typiquement les visiteurs non logués --- ne contribue à aucune arête. \\
\addlinespace
H9 & \cell{3} & \code{emitted.between(start\_date, end\_date)} : les bornes sont des \emph{dates} converties en timestamps à minuit, donc la journée du \code{end\_date} est de fait exclue. & La fenêtre réelle est de 120 jours et non 121 ; un décalage d'un jour entre le graphe et un autre calcul qui utiliserait la même variable. \\
\addlinespace
H10 & \code{sessionize\_views} & Une session = 30 minutes d'inactivité, par utilisateur réconcilié, tous appareils confondus. Le tri \code{(emitted, internalId)} rend la coupure déterministe en cas d'\emph{ex æquo}. & Un temps de réflexion long scinde une intention d'achat en deux sessions ; deux intentions rapprochées fusionnent. Voir aussi \hyp{V12} pour la sémantique du \code{cast("long")}. \\
\addlinespace
H11 & \cell{4} & Le filtre catalogue s'applique \emph{avant} la sessionisation. & Retirer une vue au milieu d'une session allonge l'écart entre ses voisines : une session peut se scinder en deux, ou perdre des produits. Le catalogue modifie donc la \emph{sessionisation}, pas seulement le jeu de nœuds. \\
\addlinespace
H12 & \code{build\_co\_view\_edges} & Deux produits vus dans la même session sont liés, indépendamment de l'ordre, de l'écart de temps et de la distance dans la session. & Une session de navigation exploratoire produit un graphe complet de produits sans rapport ; le signal séquentiel (A puis B) est perdu. \\
\addlinespace
H13 & \code{build\_co\_view\_edges} & Une session compte une fois par paire (\code{distinct} avant la self-jointure) : le poids d'une arête = nombre de sessions distinctes contenant la paire. & Un utilisateur qui compare intensément deux produits ne pèse pas plus qu'une co-vue fortuite au sein d'une session. \\
\addlinespace
H14 & \code{build\_co\_view\_edges} & \code{USE\_SESSION\_WEIGHT\_NORM = False} : toutes les sessions pèsent pareil. Une session de $n$ produits produit $\binom{n}{2}$ paires de poids 1. & Une session de 50 produits (bot, crawler, comparateur) apporte 1225 arêtes contre 1 pour une session de 2 produits : le graphe est dominé par les sessions longues. Dans cette configuration \code{weight} et \code{raw\_count} sont identiques. \\
\addlinespace
H15 & \code{build\_co\_view\_edges} & Aucun seuil minimal de poids : une co-occurrence unique est une arête à part entière. & Le graphe est majoritairement constitué d'arêtes de poids 1, c'est-à-dire de coïncidences --- et ce sont exactement celles qui varient d'un run à l'autre (\hyp{V6}). \\
\addlinespace
H16 & \code{build\_co\_view\_edges}, \code{edges\_to\_tensors} & La co-vue est symétrique et sans boucle : \code{A < B} canonise la paire, puis \code{symmetric=True} la duplique dans les deux sens. & Aucune asymétrie de type « B est un accessoire de A » n'est représentable. \\
\addlinespace
H17 & \code{split\_sessions} & Découpage temporel à la granularité \emph{session} (\code{session\_start}), pas à celle de l'événement. & Une session commencée avant \code{cutoff\_train} et poursuivie après reste entièrement dans train : de l'information postérieure à la coupure entre dans l'entraînement. \\
\addlinespace
H18 & \code{remove\_train\_edges} & Une paire évaluée doit être inédite vis-à-vis du graphe d'entraînement. & Sans cela on mesurerait la mémorisation. Noter que val et test ne sont pas dédoublonnés \emph{entre eux} : une même paire peut être dans les deux. \\
\addlinespace
H19 & \cell{5}, \code{build\_node\_mapping} & Le jeu de nœuds est l'union train $\cup$ val $\cup$ test. & Des nœuds existent dans le graphe sans aucune arête d'entraînement (isolés), uniquement parce qu'ils ont été co-vus pendant val ou test. Ils augmentent \code{num\_nodes} et donc la taille des couches d'embedding. Les tenseurs \code{val\_edge\_index} / \code{test\_edge\_index} de co-vue ne servent, eux, qu'aux compteurs. \\
\addlinespace
H20 & \code{build\_node\_mapping} & Les index de nœuds sont l'ordre croissant des \code{internalId}. & Reproductible à jeu de nœuds identique ; mais l'ajout d'un seul produit décale tous les index suivants et rend un \code{data.pt} incompatible avec un \code{node2idx.json} d'un autre run (\hyp{V13}). \\

\end{longtable}
}
\refstepcounter{table}\label{tab:hyp-a}
\begin{center}\footnotesize\itshape Tableau \thetable{} --- Hypothèses de la chaîne A.\end{center}

%======================================================================
\clearpage
\section{Hypothèses de la chaîne B --- les arêtes d'entraînement}
\label{sec:hyp-b}

{\small
\begin{longtable}{@{}l p{2.45cm} p{5.9cm} p{4.5cm}@{}}
\toprule
\textbf{Id} & \textbf{Étape} & \textbf{Hypothèse métier} & \textbf{Si elle est fausse} \\
\midrule
\endfirsthead
\toprule
\textbf{Id} & \textbf{Étape} & \textbf{Hypothèse métier} & \textbf{Si elle est fausse} \\
\midrule
\endhead
\bottomrule
\endfoot

H21 & \code{recommended\_products\_for\_t2s\_user} & Un log ad et une vue produit du même cookie, du même produit et du même client, séparés de moins de 3 secondes, sont le \emph{même} événement. & Fenêtre trop large : plusieurs vues apparient un même log et dupliquent les lignes. Trop étroite : un log est perdu par latence réseau ou décalage d'horloge. C'est le paramètre le plus sensible de la chaîne~B. \\
\addlinespace
H22 & idem & Seules les pages \code{PRODUCT} portent une recommandation exploitable. & Les recommandations de la page d'accueil, de la recherche et des pages catégorie sont hors périmètre. \\
\addlinespace
H23 & \cell{11} & Le triplet exact (\code{internalId}, \code{emitted}, \code{userId}) identifie la vue déclencheuse, donc rattache l'exécution à sa session. & Un double envoi du tracking au même timestamp duplique l'exécution. Et si la vue déclencheuse n'a pas survécu au filtre catalogue (\hyp{H2}), l'exécution est perdue (\hyp{E11}). \\
\addlinespace
H24 & \cell{12}, \cell{18} & L'ordre du tableau \code{relevantProducts} est l'ordre d'affichage réel, et seuls les 12 premiers comptent. & Le commentaire du code prévoit explicitement la variante par score (\code{relevanceScore}, \code{performanceScore}, \code{fusionScore}) ; les cellules 17--23 la calculent et mesurent l'écart entre les deux ordres. \\
\addlinespace
H25 & \cell{12}, \cell{18} & Tous les placements sponsorisés d'une exécution sont fusionnés, le rang retenu est le rang minimal atteint. Le \code{dropDuplicates(["executionId","placement\_pos"])} suppose que les lignes dupliquées par la jointure session portent des \code{candidates} identiques. & Deux placements concurrents sur une même page sont traités comme une seule liste ordonnée. \\
\addlinespace
H26 & \cell{12}, \cell{15} & \textbf{Positif} = candidat recommandé \emph{et} vu dans la même session --- sans contrainte temporelle : un produit vu \emph{avant} la recommandation compte comme positif. & C'est une fuite temporelle dans la définition du label : une partie des « clics prédits » sont des produits déjà visités, que le modèle peut retrouver par le graphe plutôt que par une vraie prédiction. \\
\addlinespace
H27 & \cell{12} & \textbf{Négatif} = candidat recommandé, non vu dans la session, \emph{et} absent des arêtes d'entraînement (testées dans les deux sens). & Évite d'apprendre contre des faux négatifs, mais retire les négatifs les plus difficiles --- ceux qui ressemblent le plus à des positifs. \\
\addlinespace
H28 & \cell{12} & Asymétrie assumée : les négatifs sont filtrés sur les arêtes du graphe d'entraînement, les positifs ne le sont pas. & Les positifs peuvent être des arêtes déjà présentes dans \code{edge\_index}, les négatifs jamais. Un modèle qui apprend seulement « cette paire est-elle une arête du graphe d'entrée ? » obtient un bon MRR sans généraliser. À quantifier : part des positifs déjà arêtes de train. \\
\addlinespace
H29 & \cell{12} & Négatifs tronqués au rang du positif le plus profond de l'exécution ; les exécutions sans positif sont écartées (\code{sample(fraction=0.0, seed=42)}). & Le modèle ne voit jamais de liste purement négative, alors que la production en produit. La \code{fraction} est un levier laissé à 0, pas une propriété du problème. \\
\addlinespace
H30 & \cell{12} & Un candidat absent de \code{node2idx} n'est pas scorable, donc \code{dropna} le supprime. & Le MRR du modèle porte sur une liste de candidats \emph{filtrée}, le MRR production affiché juste au-dessus sur la liste \emph{complète} : les deux ne sont pas strictement comparables tant que le taux de perte n'est pas mesuré. \\
\addlinespace
H31 & \cell{12} & \code{exec2code} couvre toutes les exécutions du split (indépendamment du label) et est rendu stable par tri. & Stable à jeu d'exécutions identique seulement : la moindre exécution en plus ou en moins renumérote tout et invalide les parquets déjà écrits (\hyp{V13}). \\
\addlinespace
H32 & \cell{20} & Pour les tables de référence production, le positif est défini par simple appartenance aux produits de la session, \emph{sans} filtrage sur les arêtes du graphe. & Choix explicite dans le notebook : on veut le vrai retour production. Conséquence : la baseline et le modèle ne sont pas évalués sur exactement la même définition de positif (\hyp{H27}). \\
\addlinespace
H33 & \cell{20} & Une exécution = un trigger : \code{dropDuplicates(["exec\_code"])}. & Si une exécution a été rattachée à plusieurs triggers ou à plusieurs sessions par la jointure $\pm 3$\,s, la ligne conservée est arbitraire, donc variable d'un run à l'autre (\hyp{V11}). \\
\addlinespace
H34 & \cell{25} & Le \emph{message passing} circule sur les co-vues, la tête de prédiction apprend sur les paires ads : deux sémantiques d'arête coexistent dans le même objet \code{Data}. & Cohérent si les co-vues sont vues comme une structure de voisinage et les paires ads comme la tâche. Il faut simplement ne jamais interpréter un score comme « probabilité de co-vue ». \\

\end{longtable}
}
\refstepcounter{table}\label{tab:hyp-b}
\begin{center}\footnotesize\itshape Tableau \thetable{} --- Hypothèses de la chaîne B.\end{center}

%======================================================================
\clearpage
\section{Ce qui peut être exclu du graphe}
\label{sec:exclusions}

Chaque filtre et chaque jointure \code{inner} du pipeline retire des lignes sans
alerte. Le tableau ci-dessous les liste dans l'ordre du flux : c'est l'inventaire à
instrumenter si l'on veut savoir combien de signal est perdu et où.

{\small
\begin{longtable}{@{}l p{3.2cm} p{4.4cm} p{5.2cm}@{}}
\toprule
\textbf{Id} & \textbf{Étape} & \textbf{Mécanisme} & \textbf{Ce qui disparaît} \\
\midrule
\endfirsthead
\toprule
\textbf{Id} & \textbf{Étape} & \textbf{Mécanisme} & \textbf{Ce qui disparaît} \\
\midrule
\endhead
\bottomrule
\endfoot

E1 & \code{get\_customer} & jointure \code{inner} clients / publishers & Un client sans \emph{publisher} \code{artemis-prod} : le notebook s'exécute alors sur des listes vides sans erreur. \\
E2 & \code{get\_products} & \code{isActive}, \code{isOkInStream} & Produits désactivés ou hors flux au moment du run (pas au moment de la vue). \\
E3 & \code{get\_products} & \code{parentId = -1} & Toutes les déclinaisons. \\
E4 & \code{get\_embeddings} & \code{last\_model\_id} figé & Produits embeddés par une autre version : pas exclus du graphe, mais réduits à un vecteur nul. \\
E5 & \code{get\_views} & \code{user.internalId} non nul & Trafic non identifié. \\
E6 & \code{get\_views} & \code{masterProductInternalId} non nul & Vues non résolues en produit maître. \\
E7 & \code{get\_views} & \code{emitted.between(...)} & Journée du \code{end\_date} (bornes prises à minuit, \hyp{H9}). \\
E8 & \cell{4} & jointure \code{inner} vues / catalogue & Vues sur des produits absents du catalogue courant --- et, par effet de bord, la sessionisation des vues restantes change (\hyp{H11}). \\
E9 & \code{build\_co\_view\_edges} & self-jointure \code{L.internalId < R.internalId} & Sessions à un seul produit distinct : elles ne produisent aucune arête, quel que soit leur volume. \\
E10 & \code{remove\_train\_edges} & \code{left\_anti} & Paires val/test déjà vues en train. \\
E11 & \code{recommended\_products\_for\_t2s\_user} & \code{pageType = "PRODUCT"}, fenêtre $\pm 3$\,s, \code{publisherId isin(...)} & Recommandations hors page produit ; logs non appariés à une vue ; logs d'un \code{publisherId} non résolu. \\
E12 & \cell{11} & jointure \code{inner} sur (\code{internalId}, \code{emitted}, \code{userId}) & Exécutions dont la vue déclencheuse n'a pas survécu à \hyp{E5}--\hyp{E8}. \\
E13 & \cell{12} & \code{explode(sponsoredProductPlacementExecutions)} & Exécutions sans placement sponsorisé (tableau vide ou nul). \\
E14 & \cell{12} & \code{slice(relevantProducts, 1, 12)} & Candidats au-delà du rang 12 --- y compris les positifs qui s'y trouvaient, ce qui plafonne mécaniquement le MRR production. \\
E15 & \cell{12} & \code{is\_existing\_edge} & Paires candidates déjà arêtes du graphe d'entraînement, côté négatifs uniquement (\hyp{H28}). \\
E16 & \cell{12} & \code{rank} $\le$ \code{deepest\_pos\_rank} & Négatifs plus profonds que le positif le plus profond. \\
E17 & \cell{12} & \code{sample(fraction=0.0, seed=42)} & \emph{Tous} les négatifs des exécutions sans positif. \\
E18 & \cell{12} & \code{map(node2idx)} + \code{dropna} & Paires dont un produit n'est pas un nœud du graphe (\hyp{H30}). \\
E19 & \cell{20} & \code{dropDuplicates(["exec\_code"])} & Triggers surnuméraires d'une même exécution, choisis arbitrairement. \\

\end{longtable}
}
\refstepcounter{table}\label{tab:exclusions}
\begin{center}\footnotesize\itshape Tableau \thetable{} --- Points d'exclusion, dans l'ordre du flux.\end{center}

\paragraph{Les deux exclusions à mesurer en priorité.}
\hyp{E8} conditionne tout l'amont : la taille du graphe est celle du catalogue au
jour du run, pas au jour des vues. \hyp{E18} conditionne toute la comparaison au
MRR production : tant que la proportion de candidats perdus n'est pas affichée, on
ne sait pas si le modèle est comparé à la baseline sur la même population.

%======================================================================
\clearpage
\section{Pourquoi deux runs donnent des graphes différents}
\label{sec:variabilite}

Le notebook ne fixe aucune version de table. Deux exécutions séparées de quelques
jours, \emph{sur la même fenêtre de dates}, lisent des données différentes. Les
causes se rangent en quatre familles.

\subsection{Sources amont non figées}

\begin{itemize}[leftmargin=*]
  \item \hyp{V1} \textbf{Le catalogue est un instantané.}
        \code{t2s\_mongo\_product\_0\_current} est, comme son nom l'indique, l'état
        \emph{courant}. \code{isActive}, \code{isOkInStream} et \code{parentId}
        évoluent : un produit désactivé entre deux runs disparaît du graphe, y
        compris pour les vues de janvier. C'est la première cause d'écart, et elle
        est monotone dans le temps --- plus on relance tard, plus le graphe rétrécit.
  \item \hyp{V2} \textbf{La table d'embeddings s'enrichit.} Le filtre sur
        \code{last\_model\_id} fige le modèle, mais pas la population : les produits
        embeddés depuis le run précédent passent du vecteur nul à un vrai vecteur.
        Le compteur \code{Nodes without embedding} en est le témoin direct.
  \item \hyp{V3} \textbf{Le tracking reçoit des données tardives.} Événements
        arrivés en retard, retraitements et réécritures de partitions modifient le
        contenu d'une fenêtre historique déjà passée.
  \item \hyp{V4} \textbf{\code{t2s\_merged\_user} évolue.} La réconciliation
        d'identités s'améliore avec le temps : deux \code{userId} distincts hier
        deviennent un seul \code{masterId} aujourd'hui. Leurs vues fusionnent, les
        sessions changent, de nouvelles paires apparaissent --- \emph{sans qu'aucun
        événement nouveau n'ait été émis}. C'est la cause la plus contre-intuitive.
  \item \hyp{V5} \textbf{Rétention des logs publicitaires.} Si
        \code{ads\_adlog\_fct} a une durée de conservation bornée, relancer plus tard
        sur la même fenêtre donne moins d'exécutions, donc moins de positifs et un
        MRR production différent. \emph{À vérifier} sur la table --- le notebook
        n'en dit rien.
\end{itemize}

\subsection{Amplification : une petite variation d'entrée, un grand écart en sortie}

Le pipeline est quadratique par construction, ce qui transforme un écart marginal
en entrée en écart massif en sortie.

\begin{itemize}[leftmargin=*]
  \item \hyp{V6} \textbf{Une vue de plus dans une session de $n$ produits ajoute
        jusqu'à $n$ arêtes.} Et comme aucun seuil de poids n'est appliqué
        (\hyp{H15}), ces arêtes de poids 1 entrent telles quelles dans le graphe.
  \item \hyp{V7} \textbf{Une fusion d'utilisateurs enchaîne deux sessions.} Si deux
        sessions de tailles $a$ et $b$ fusionnent, on passe de
        $\binom{a}{2} + \binom{b}{2}$ à $\binom{a+b}{2}$ paires, soit
        $a \times b$ arêtes créées d'un coup.
  \item \hyp{V8} \textbf{Retirer une vue intermédiaire peut scinder une session en
        deux} (\hyp{H11}) : les paires qui traversaient la coupure disparaissent
        toutes ensemble.
  \item \hyp{V9} \textbf{Les sessions longues dominent} (\hyp{H14}) : une seule
        session de bot en plus ou en moins déplace des milliers d'arêtes.
\end{itemize}

\subsection{Non-déterminisme interne au notebook (mêmes entrées, sortie différente)}

\begin{itemize}[leftmargin=*]
  \item \hyp{V10} \code{node\_features[positions] = emb\_matrix} (\cell{8}) : si un
        \code{internalId} a plusieurs lignes d'embedding, la dernière écrite gagne,
        dans l'ordre de collecte de Spark --- qui n'est pas garanti.
  \item \hyp{V11} \code{dropDuplicates(["exec\_code"])} (\cell{20}) et
        \code{dropDuplicates(["executionId","placement\_pos"])} (\cell{18}) :
        Spark conserve une ligne \emph{arbitraire}, pas une ligne définie par un tri.
  \item \hyp{V12} \code{(emitted - lag(emitted)).cast("long")}
        (\code{sessionize\_views}) : le résultat de la soustraction de deux
        timestamps, et sa conversion en \code{long}, dépendent de la version de Spark
        et du mode d'intervalles ANSI. Le code suppose des \emph{secondes}, cohérent
        avec \code{session\_timeout = 1800}. \emph{À vérifier} avant tout changement
        d'environnement : une autre unité ferait de chaque vue une session isolée, et
        le graphe s'effondrerait silencieusement.
  \item \hyp{V13} \textbf{Couplage des artefacts.} \code{node2idx} (\hyp{H20}) et
        \code{exec\_code} (\hyp{H31}) sont deux renumérotations globales : elles sont
        stables à jeu constant, et entièrement décalées dès qu'un nœud ou une
        exécution apparaît. Un \code{data.pt} ne peut donc être lu qu'avec le
        \code{node2idx.json}, les \code{exec\_mappings.json} et les parquets
        \emph{du même run}. C'est exactement ce que dit le commentaire de la
        \cell{28} pour justifier le suffixe \code{\_prototype}.
  \item \hyp{V14} \textbf{Un run long peut chevaucher une mise à jour amont.}
        \code{df\_split} est mis en cache, mais \code{df\_train\_edges} ne l'est pas
        et il est réévalué à chaque usage (compteurs, \code{df\_graph\_nodes},
        \code{toPandas}, puis trois appels à
        \code{positive\_negative\_edge\_construction}). Si le cache est évincé ou si
        une table amont change en cours de run, deux parties du même graphe peuvent
        être calculées sur deux instantanés différents.
\end{itemize}

\subsection{Paramètres qui changent le graphe à données constantes}

\begin{center}\small
\begin{tabular}{@{}l l p{7.6cm}@{}}
\toprule
\textbf{Paramètre} & \textbf{Valeur} & \textbf{Effet} \\
\midrule
\code{session\_timeout} & \code{1800} & Frontières de session, donc toutes les paires (\hyp{H10}). \\
\code{start\_date} / \code{end\_date} & \code{2026-01-01} / \code{2026-05-01} & Volume total ; dernier jour exclu (\hyp{H9}). \\
\code{cutoff\_train} / \code{cutoff\_val} & $-30$\,j / $-15$\,j & Répartition train/val/test (\hyp{H17}). \\
\code{USE\_SESSION\_WEIGHT\_NORM} & \code{False} & Poids des sessions longues (\hyp{H14}). \\
\code{last\_model\_id} & \code{06bbe57\dots} & Population des nœuds avec features (\hyp{H5}). \\
fenêtre adlog/tracking & $\pm 3$\,s & Nombre d'exécutions appariées (\hyp{H21}). \\
\code{slice(..., 1, 12)} / \code{TOP\_K} & \code{12} & Plafond du MRR production (\hyp{E14}). \\
\code{sample(fraction=...)} & \code{0.0} & Négatifs des exécutions sans positif (\hyp{H29}). \\
\code{symmetric} & \code{True} & Nombre d'arêtes doublé (\hyp{H16}). \\
seuil de poids minimal & \emph{absent} & Bruit et instabilité du graphe (\hyp{H15}). \\
\bottomrule
\end{tabular}
\end{center}

\subsection{Rendre un run reproductible}

Par ordre d'efficacité, et sans changer la logique métier :

\begin{enumerate}[leftmargin=*]
  \item \textbf{Figer les instantanés.} Lire chaque table en \emph{time travel} Delta
        (\code{VERSION AS OF} / \code{TIMESTAMP AS OF}) et journaliser la version
        retenue à côté du \code{data.pt}. C'est ce qui neutralise
        \hyp{V1}--\hyp{V5} et \hyp{V14} d'un seul coup.
  \item \textbf{Reconstruire le catalogue à la date de la vue} plutôt qu'au jour du
        run, si une table historisée existe --- sinon assumer \hyp{H2} et le noter
        dans les métadonnées du dataset.
  \item \textbf{Rendre les déduplications déterministes} : remplacer les
        \code{dropDuplicates} de \hyp{V11} par un \code{row\_number} sur un tri
        explicite, et dédoublonner les embeddings sur
        \code{max(last\_update\_date)} avant la \cell{8} (\hyp{V10}).
  \item \textbf{Publier un manifeste de run} : versions de tables, valeurs des
        paramètres du tableau ci-dessus, \code{num\_nodes}, nombre d'arêtes par
        split, \code{Nodes without embedding}, taux de perte de \hyp{E18}. Deux runs
        deviennent alors comparables même s'ils diffèrent.
  \item \textbf{Instrumenter les exclusions} \hyp{E1}--\hyp{E19} par un compteur de
        lignes avant/après. Un écart inexpliqué entre deux runs se localise ainsi en
        une lecture.
  \item \textbf{Introduire un seuil de poids minimal} (\hyp{H15}) et/ou une
        normalisation par taille de session (\hyp{H14}) : c'est le levier le plus
        direct pour stabiliser le graphe, puisqu'il supprime précisément les arêtes
        qui varient.
\end{enumerate}

%======================================================================
\section{Les cinq points à trancher}

Résumé pour décision, hors reproductibilité :

\begin{enumerate}[leftmargin=*]
  \item \hyp{H26} --- un positif peut être un produit vu \emph{avant} la
        recommandation. À restreindre à \code{emitted > logInstant} si l'on veut
        mesurer une prédiction et non une co-occurrence.
  \item \hyp{H28} --- les négatifs excluent les arêtes du graphe, les positifs non.
        Mesurer la part de positifs déjà présents dans \code{edge\_index} : c'est le
        plafond de ce que le modèle peut obtenir sans généraliser.
  \item \hyp{H30} et \hyp{E18} --- MRR modèle sur candidats filtrés contre MRR
        production sur candidats complets. Afficher le taux de perte, sinon les deux
        chiffres ne se comparent pas.
  \item \hyp{H14} et \hyp{H15} --- absence de normalisation par session et de seuil
        de poids : le graphe est dominé par les sessions longues et les coïncidences.
  \item \hyp{H19} --- les nœuds val/test sans arête d'entraînement gonflent
        \code{num\_nodes}. Décider s'ils doivent être des nœuds isolés du graphe
        d'entraînement ou en être retirés.
\end{enumerate}

\end{document}
